\chapter{L'apparato muscolo-scheletrico}

% ---------------------------------------------------------------------------
\section{Generalità sull'apparato muscolare}

\textbf{Muscolo} (lat.\ \textit{musculus}): organo di \textbf{tessuto muscolare} (contrattile);
delimita le regioni ed è effettore dei movimenti. In un muscolo isolato: \textbf{ventre} (porzione
carnosa centrale), \textbf{testa} e \textbf{coda} (estremità), raccordate all'osso dai
\textbf{tendini}. Può essere normale o atrofizzato (ventre ridotto).

\subsection{Classificazione funzionale}
\begin{itemize}
  \item \textbf{agonisti}: compiono il movimento;
  \item \textbf{antagonisti}: si oppongono agli agonisti (movimenti opposti, rallentamento/arresto);
  \item \textbf{sinergici}: concorrono allo stesso movimento (aumentano forza, velocità, modulano
    l'intensità).
\end{itemize}

% ---------------------------------------------------------------------------
\section{Muscoli della testa e del collo}

Permettono i movimenti di faccia, lingua, faringe, laringe (comunicazione, masticazione,
deglutizione). Due gruppi: \textbf{masticatori} (movimenti della mandibola) e \textbf{mimici /
pellicciai} (tegumenti facciali).

\subsection{Muscoli masticatori}
Tutti pari.
\begin{itemize}
  \item \textbf{temporale}: linea temporale inferiore/squama del temporale $\rightarrow$ processo
    coronoideo; innalza la mandibola, chiude la bocca, retrazione e movimenti laterali;
  \item \textbf{massetere}: arco zigomatico $\rightarrow$ angolo mandibolare; innalza, chiude,
    protrusione, retrazione, movimenti laterali;
  \item \textbf{pterigoideo laterale}: grande ala dello sfenoide e processo pterigoideo $\rightarrow$
    condilo mandibolare e ATM; apertura della bocca, protrusione, movimenti laterali;
  \item \textbf{pterigoideo mediale}: processo pterigoideo, palatino, tuberosità della mascella
    $\rightarrow$ angolo mandibolare; innalza, chiude, movimenti laterali.
\end{itemize}

\subsection{Muscoli mimici}
Innervati dal \textbf{nervo facciale}; espressione facciale (sopracciglia, palpebre, naso, bocca).
\begin{itemize}
  \item \textbf{temporo-parietale}: fascia temporale $\rightarrow$ galea aponeurotica;
  \item \textbf{estrinseci del padiglione auricolare}: galea aponeurotica $\rightarrow$ padiglione;
  \item \textbf{nasale}: parte trasversa (gioghi alveolari di canini/incisivi superiori $\rightarrow$
    margine della narice) e parte alare (apertura piriforme $\rightarrow$ dorso e ala del naso);
  \item altri: ventre frontale dell'occipito-frontale, corrugatore del sopracciglio, orbicolare
    dell'occhio, zigomatico minore/maggiore, elevatore del labbro superiore, procero, orbicolare
    della bocca, risorio, platisma, mentale, depressore dell'angolo della bocca, depressore del
    labbro inferiore.
\end{itemize}

\subsection{Muscoli estrinseci dell'occhio}
Governano il bulbo oculare:
\begin{itemize}
  \item \textbf{retto inferiore}: guardare in basso; \textbf{retto mediale}: rotazione mediale;
  \item \textbf{retto superiore}: guardare in alto; \textbf{retto laterale}: rotazione laterale;
  \item \textbf{obliquo inferiore}: in alto e lateralmente; \textbf{obliquo superiore}: in basso e
    lateralmente;
  \item + elevatore della palpebra superiore; nervi: oculomotore (III), trocleare (IV), abducente
    (VI), ottico (II).
\end{itemize}

\subsection{Muscoli della lingua (estrinseci)}
\begin{itemize}
  \item \textbf{genioglosso}: abbassamento e protrusione;
  \item \textbf{ioglosso}: abbassamento e retrazione;
  \item \textbf{palatoglosso}: innalza la lingua, abbassa il palato molle;
  \item \textbf{stiloglosso}: retrazione, innalzamento delle parti laterali.
\end{itemize}
In rapporto con mandibola, osso ioide, processo stiloideo del temporale.

\subsection{Muscoli anteriori del collo}
\rubrica{Sopraioidei}
\begin{itemize}
  \item \textbf{digastrico}: ventre anteriore (fossetta digastrica della mandibola) + posteriore
    (incisura mastoidea) $\rightarrow$ osso ioide (tendine intermedio); innalza lo ioide, abbassa la
    mandibola; innervato da nervo mandibolare (V, ventre ant.) e facciale (VII, ventre post.);
  \item \textbf{stiloioideo}: processo stiloideo $\rightarrow$ corpo dello ioide; lo innalza e
    retrae (nervo facciale);
  \item \textbf{miloioideo}: linea miloioidea della mandibola $\rightarrow$ corpo dello ioide e rafe;
    innalza lo ioide, solleva il pavimento orale e la lingua, abbassa la mandibola (nervo
    mandibolare);
  \item \textbf{genioioideo}: spina mentale della mandibola $\rightarrow$ corpo dello ioide; innalza
    lo ioide (con la lingua in avanti), abbassa la mandibola (C1 via ipoglosso, XII).
\end{itemize}
\rubrica{Sottoioidei}
\begin{itemize}
  \item \textbf{omoioideo}: margine superiore della scapola $\rightarrow$ corpo/grande corno dello
    ioide; lo abbassa e retrae (C1--C3);
  \item \textbf{sternoioideo}: clavicola/manubrio $\rightarrow$ corpo dello ioide; lo abbassa
    (C1--C3);
  \item \textbf{sternotiroideo}: manubrio e I cartilagine costale $\rightarrow$ cartilagine tiroidea;
    abbassa la tiroidea e lo ioide (C1--C3);
  \item \textbf{tiroioideo}: cartilagine tiroidea $\rightarrow$ corpo/grande corno dello ioide;
    innalza la tiroidea, abbassa lo ioide (C1 via ipoglosso).
\end{itemize}

\subsection{Muscolatura laterale del collo}
\begin{itemize}
  \item \textbf{sternocleidomastoideo}: bipennato, 4 capi; sterno e clavicola $\rightarrow$ processo
    mastoideo e linea nucale superiore; rotazione controlaterale, inclinazione omolaterale,
    estensione; inspiratorio (solleva sterno e clavicola); nervo accessorio;
  \item \textbf{platisma}: II costa/pettorali e clavicola $\rightarrow$ regione masseterina,
    mandibola, commessura labiale; abbassa mandibola e commessura labiale.
\end{itemize}

\subsection{Muscoli prevertebrali}
Flettono e inclinano lateralmente il capo:
\begin{itemize}
  \item \textbf{retto anteriore della testa}: atlante $\rightarrow$ occipitale;
  \item \textbf{lungo del collo}: prime 3 toraciche/ultime 3 cervicali $\rightarrow$ prime cervicali;
  \item \textbf{lungo della testa}: occipitale $\rightarrow$ vertebre cervicali.
\end{itemize}

\subsection{Muscoli suboccipitali}
Estensori e rotatori del capo; 4, pari, profondi (articolazioni cranio-vertebrali):
\begin{itemize}
  \item \textbf{piccolo retto posteriore della testa}: tubercolo posteriore dell'atlante
    $\rightarrow$ occipitale;
  \item \textbf{grande retto posteriore della testa}: processo spinoso dell'epistrofeo $\rightarrow$
    occipitale;
  \item \textbf{obliquo superiore}: processo trasverso dell'atlante $\rightarrow$ occipitale (presso
    il mastoideo);
  \item \textbf{obliquo inferiore}: processo spinoso dell'epistrofeo $\rightarrow$ processo trasverso
    dell'atlante.
\end{itemize}

% ---------------------------------------------------------------------------
\section{Muscoli del torace}

\begin{itemize}
  \item \textbf{estrinseci} (torace $\rightarrow$ cinture/colonna): gran pettorale, piccolo
    pettorale, succlavio, dentato anteriore, diaframma;
  \item \textbf{intrinseci} (nel torace): intercostali (interni/esterni/intimi), trasverso del
    torace, elevatori delle coste, sottocostali.
\end{itemize}

\subsection{Gran pettorale}
Tre parti: \textbf{clavicolare} (metà mediale della clavicola), \textbf{sternocostale} (sterno e
II--VI cartilagine costale), \textbf{addominale} (linea alba) $\rightarrow$ cresta del tubercolo
maggiore dell'omero. Adduce e intraruota l'omero; a punto fisso sull'omero solleva il tronco;
inspiratorio accessorio.

\subsection{Piccolo pettorale}
III--V costa $\rightarrow$ processo coracoideo. A punto fisso sul torace abbassa la scapola; a punto
fisso sulla scapola solleva le coste (inspiratorio).

\subsection{Succlavio}
I costa $\rightarrow$ faccia inferiore della clavicola. Abbassa la spalla; a punto fisso sulla
clavicola innalza la costa (inspiratorio); protegge plesso brachiale e vasi succlavi.

\subsection{Dentato anteriore}
Prime 10 coste $\rightarrow$ margine mediale della scapola. Sposta in avanti la scapola; a punto
fisso sulla scapola eleva le coste (inspiratorio).

\subsection{Diaframma}
Lamina muscolo-fibrosa (doppia cupola) che separa torace e addome; principale muscolo respiratorio.
Origini tripartite (sternale, costale, lombare) $\rightarrow$ centro tendineo.
\begin{itemize}
  \item \textbf{inspirazione}: discesa del centro tendineo (favorita dal rilassamento addominale)
    $\rightarrow$ poi innalzamento delle coste e ampliamento della gabbia;
  \item \textbf{espirazione}: intercostali interni + diaframma; se forzata, anche muscoli addominali
    (retto, obliquo esterno/interno, trasverso) $\rightarrow$ $\uparrow$ pressione intraddominale
    $\rightarrow$ solleva il centro tendineo e abbassa le coste (soffiarsi il naso, tosse, starnuto).
\end{itemize}

\subsection{Elevatori delle coste e sottocostali}
\begin{itemize}
  \item \textbf{elevatori delle coste}: processi trasversi di C7 e T9--T11 $\rightarrow$ costa
    sottostante; inspiratori;
  \item \textbf{sottocostali}: faccia interna delle coste $\rightarrow$ costa sottostante/successiva;
    espiratori.
\end{itemize}

\subsection{Muscoli intercostali}
\begin{itemize}
  \item \textbf{esterni}: margine inferiore della costa $\rightarrow$ margine superiore della
    sottostante; elevano le coste (inspirazione);
  \item \textbf{interni}: labbro mediale del solco costale $\rightarrow$ margine superiore della
    costa/cartilagine sottostante; abbassano le coste (espirazione);
  \item \textbf{intimi}: margine inferiore $\rightarrow$ margine superiore della sottostante.
\end{itemize}

% ---------------------------------------------------------------------------
\section{Muscoli dell'addome}

\begin{itemize}
  \item \textbf{anterolaterali}: retto, piramidale, obliquo interno/esterno, trasverso
    $\rightarrow$ tono = pressione endoaddominale (mantiene gli organi in sede); la contrazione ne
    favorisce lo svuotamento;
  \item \textbf{posteriori}: quadrato dei lombi, ileopsoas.
\end{itemize}

\subsection{Retto dell'addome}
Processo xifoideo e V--VII cartilagine costale $\rightarrow$ ramo superiore del pube. Abbassa le
coste (espirazione), flette torace/pelvi, aumenta la pressione addominale (il ``\textit{six pack}'').

\subsection{Obliquo esterno}
Faccia esterna delle ultime 8 coste $\rightarrow$ cresta iliaca e aponeurosi; l'aponeurosi forma la
\textbf{linea alba}. Abbassa le coste; asimmetrico $\rightarrow$ rotazione controlaterale del tronco;
sinergico $\rightarrow$ flette torace su pelvi.

\subsection{Obliquo interno}
Legamento inguinale, spina iliaca ant.\ sup., cresta iliaca, fascia toracolombare $\rightarrow$
ultime 3 cartilagini costali e aponeurosi. Ruota il torace omolateralmente; bilaterale $\rightarrow$
flette la colonna, abbassa le coste, aumenta la pressione addominale.

\subsection{Piramidale}
Ramo superiore del pube $\rightarrow$ linea alba (la tende); assente in $\sim$1/6 dei soggetti.

\subsection{Trasverso dell'addome}
Il più profondo. Ultime 6 cartilagini costali, fascia toracolombare, cresta iliaca, legamento
inguinale $\rightarrow$ aponeurosi. Abbassa le coste, aumenta la pressione addominale; il tempismo
di trasverso/obliqui/retto stabilizza la colonna lombare e previene la lombalgia.

\subsection{Quadrato dei lombi}
Processi costiformi L2--L4 e cresta iliaca $\rightarrow$ processi costiformi L1--L4 e XII costa.
Abbassa la XII costa, inclina lateralmente colonna lombare e pelvi.

\subsection{Grande psoas e iliaco}
\begin{itemize}
  \item \textbf{grande psoas}: processi costiformi L1--L5 e corpi/dischi T12--L5 $\rightarrow$
    piccolo trocantere; flette e ruota lateralmente la coscia; a femore fisso flette tronco e bacino;
  \item \textbf{iliaco}: fossa iliaca $\rightarrow$ piccolo trocantere; flessore che completa lo
    psoas $\rightarrow$ insieme = \textbf{ileopsoas}.
\end{itemize}

% ---------------------------------------------------------------------------
\section{Muscoli del dorso}

Sostengono la postura contro gravità. Tre gruppi:
\begin{itemize}
  \item \textbf{spino-appendicolari}: trapezio, gran dorsale, romboidi, elevatore della scapola;
  \item \textbf{spino-costali}: dentato posteriore superiore e inferiore;
  \item \textbf{spino-dorsali}: spino-trasversari (splenio), erettore della colonna,
    trasverso-spinali, interspinosi, intertrasversari.
\end{itemize}

\subsection{Spino-appendicolari}
\begin{itemize}
  \item \textbf{trapezio}: discendente (protuberanza occipitale e legamento nucale $\rightarrow$
    clavicola), trasversa (C7--T3 $\rightarrow$ acromion e spina della scapola), ascendente (T3--T12
    $\rightarrow$ spina della scapola); eleva e adduce la scapola; a punto fisso estende/inclina la
    testa;
  \item \textbf{gran dorsale}: T6--T12 e vertebre lombari $\rightarrow$ tubercolo minore dell'omero;
    adduce, estende, intraruota l'omero; a punto fisso eleva tronco e coste (arrampicata);
  \item \textbf{piccolo romboide}: legamento nucale e C7 $\rightarrow$ scapola; \textbf{grande
    romboide}: T1--T4 $\rightarrow$ scapola; entrambi la spostano medialmente;
  \item \textbf{elevatore della scapola}: processi trasversi C1--C4 $\rightarrow$ angolo superiore
    della scapola; la solleva e sposta medialmente.
\end{itemize}

\subsection{Spino-costali}
\begin{itemize}
  \item \textbf{dentato posteriore superiore}: C6--C7 e T1--T2 $\rightarrow$ II--V costa; le eleva;
  \item \textbf{dentato posteriore inferiore}: T11--T12 e L1--L2 $\rightarrow$ ultime 4 coste; le
    abbassa;
  \item entrambi partecipano ai movimenti respiratori.
\end{itemize}

\subsection{Spino-dorsali}
\begin{itemize}
  \item \textbf{splenio}: della testa (legamento nucale, C7--T2 $\rightarrow$ mastoideo; inclina e
    ruota omolateralmente, bilaterale estende) e del collo (estende la colonna cervicale);
  \item \textbf{erettore della colonna} (ileocostale, lunghissimo; porzioni cervicale/toracica/
    lombare): estende la colonna (stazione eretta); unilaterale la inclina;
  \item \textbf{trasverso-spinali} (semispinali, rotatori, multifido): bilaterale estende/stabilizza,
    unilaterale inclina omolateralmente e ruota controlateralmente;
  \item \textbf{interspinosi}: estensori; \textbf{intertrasversari}: inclinano omolateralmente,
    bilaterale stabilizzano.
\end{itemize}

% ---------------------------------------------------------------------------
\section{Muscoli dell'arto superiore}

\subsection{Muscoli della spalla}
La spalla = 5 articolazioni (3 vere/diartrosi, 2 spurie). Muscoli: deltoide + cuffia dei rotatori
(sovraspinato, infraspinato, piccolo rotondo, grande rotondo, sottoscapolare).
\begin{itemize}
  \item \textbf{deltoide}: 3 capi (clavicolare, acromiale, spinoso) $\rightarrow$ tuberosità
    deltoidea; anteriore = flessione/adduzione/intrarotazione, posteriore = estensione/abduzione/
    extrarotazione, medio = massima abduzione;
  \item \textbf{cuffia dei rotatori}: stabilizza la scapolo-omerale (coatta la testa dell'omero
    nella glenoide):
    \begin{itemize}
      \item \textbf{sottoscapolare}: fossa sottoscapolare $\rightarrow$ tubercolo minore;
        adduce/intraruota;
      \item \textbf{sovraspinato}: fossa sopraspinata $\rightarrow$ tubercolo maggiore;
        abduce/extraruota;
      \item \textbf{infraspinato}: fossa infraspinata $\rightarrow$ tubercolo maggiore; extraruota;
      \item \textbf{piccolo rotondo}: margine laterale della scapola $\rightarrow$ tubercolo
        maggiore; extraruota/adduce/estende;
      \item \textbf{grande rotondo}: angolo inferiore della scapola $\rightarrow$ tubercolo minore;
        adduce/intraruota/estende.
    \end{itemize}
\end{itemize}

\subsection{Muscoli anteriori del braccio}
\begin{itemize}
  \item \textbf{bicipite brachiale}: capo lungo (tubercolo sopraglenoideo) + capo breve (processo
    coracoideo) $\rightarrow$ tuberosità del radio; stabilizza la spalla, flette l'avambraccio,
    supinatore (avambraccio prono);
  \item \textbf{brachiale}: faccia anteriore dell'omero $\rightarrow$ tuberosità dell'ulna; flette
    l'avambraccio;
  \item \textbf{coracobrachiale}: processo coracoideo $\rightarrow$ faccia anteromediale dell'omero;
    flette e adduce l'omero.
\end{itemize}

\subsection{Muscoli posteriori del braccio}
\begin{itemize}
  \item \textbf{tricipite brachiale}: capo lungo (tubercolo sottoglenoideo), laterale e mediale
    (faccia posteriore dell'omero) $\rightarrow$ olecrano; adduce/estende l'omero, estende
    l'avambraccio;
  \item \textbf{anconeo}: epicondilo laterale $\rightarrow$ margine laterale dell'olecrano; estende
    l'avambraccio.
\end{itemize}

\subsection{Cavità ascellare}
Spazio piramidale tra torace e arto superiore. Pareti:
\begin{itemize}
  \item \textbf{anteriore}: gran pettorale (+ succlavio, piccolo pettorale);
  \item \textbf{posteriore}: sottoscapolare, grande rotondo, gran dorsale;
  \item \textbf{laterale}: capo breve del bicipite, coracobrachiale;
  \item \textbf{mediale}: dentato anteriore, parete toracica;
  \item \textbf{base}: fascia ascellare (linfonodi ascellari nel tessuto adiposo); \textbf{apice}:
    tra clavicola, I costa e processo coracoideo (vasi succlavi/ascellari, plesso brachiale).
\end{itemize}

\subsection{Muscoli dell'avambraccio}
Tre gruppi: \textbf{anteriori} (flessori/pronatori, epitrocleari), \textbf{laterali}, \textbf{posteriori}
(estensori/supinatori, epicondiloidei).
\rubrica{Anteriori (4 strati)}
\begin{itemize}
  \item \textbf{pronatore rotondo} (epicondilo mediale + coronoideo $\rightarrow$ faccia laterale del
    radio): prona e flette l'avambraccio;
  \item \textbf{flessore radiale del carpo} (epicondilo mediale $\rightarrow$ II metacarpale): flette
    e adduce la mano;
  \item \textbf{palmare lungo} (epicondilo mediale $\rightarrow$ aponeurosi palmare): flette la mano;
  \item \textbf{flessore ulnare del carpo} (epicondilo mediale + olecrano $\rightarrow$ pisiforme):
    flette e adduce la mano;
  \item \textbf{flessore superficiale delle dita} (epicondilo mediale/coronoideo/radio $\rightarrow$
    falangi medie II--V): le flette;
  \item \textbf{flessore profondo delle dita} (ulna/membrana interossea $\rightarrow$ falangi
    distali II--V): le flette;
  \item \textbf{flessore lungo del pollice} (radio/membrana interossea $\rightarrow$ falange distale
    del pollice): la flette;
  \item \textbf{pronatore quadrato} (ulna $\rightarrow$ radio): prona l'avambraccio.
\end{itemize}
\rubrica{Posteriori}
\begin{itemize}
  \item superficiali: \textbf{estensore delle dita} (epicondilo laterale $\rightarrow$ falangi II/III
    di II--V; le estende), \textbf{estensore del mignolo}, \textbf{estensore ulnare del carpo}
    ($\rightarrow$ V metacarpale; estende e abduce la mano);
  \item profondi: \textbf{supinatore} (epicondilo $\rightarrow$ radio; supinazione completa),
    \textbf{abduttore lungo del pollice} ($\rightarrow$ I metacarpale), \textbf{estensore breve/lungo
    del pollice}, \textbf{estensore dell'indice}.
\end{itemize}
\rubrica{Laterali}
\begin{itemize}
  \item \textbf{brachioradiale} (omero $\rightarrow$ processo stiloideo del radio): flette
    l'avambraccio (posizione intermedia);
  \item \textbf{estensore radiale lungo del carpo} ($\rightarrow$ II metacarpale) e \textbf{breve}
    ($\rightarrow$ III metacarpale): estendono e abducono la mano.
\end{itemize}

\subsection{Muscoli della mano}
Tutti palmari, tre gruppi:
\begin{itemize}
  \item \textbf{eminenza tenar} (laterali): abduttore breve, flessore breve, opponente e adduttore
    del pollice;
  \item \textbf{eminenza ipotenar} (mediali): palmare breve, abduttore, flessore breve e opponente
    del mignolo;
  \item \textbf{palmari} (intermedi): lombricali (tendini del flessore profondo $\rightarrow$
    estensore delle dita), interossei palmari e dorsali.
\end{itemize}

% ---------------------------------------------------------------------------
\section{Muscoli dell'arto inferiore}

Arto inferiore = locomozione; 6 regioni, 3 articolazioni principali: \textbf{anca} (coxo-femorale),
\textbf{coscia}, \textbf{ginocchio} (femore--tibia, rotula), \textbf{gamba} (polpaccio/stinco),
\textbf{caviglia} (tibio-tarsica, malleoli, tendine d'Achille), \textbf{piede}.

\subsection{Muscoli dell'anca}
Interni (ileo-psoas, piccolo psoas) ed esterni (glutei, gemelli, otturatori, piriforme, quadrato del
femore, tensore della fascia lata).
\begin{itemize}
  \item \textbf{grande psoas} (T12, L1--L4) + \textbf{iliaco} (fossa iliaca) = \textbf{ileo-psoas}
    $\rightarrow$ piccolo trocantere; a punto fisso sulla colonna flette/adduce/intraruota la coscia,
    a punto fisso sul femore flette il tronco;
  \item \textbf{piccolo psoas}: ultima toracica/L1 $\rightarrow$ eminenza ileopubica; flessione del
    tronco;
  \item \textbf{tensore della fascia lata}: spina iliaca ant.\ sup.\ $\rightarrow$ condilo laterale
    della tibia (tratto ileotibiale); tende la fascia lata, flette e abduce la coscia, stabilizza il
    ginocchio;
  \item \textbf{grande gluteo}: cresta iliaca, sacro, coccige $\rightarrow$ tuberosità glutea;
    estende ed extraruota la coscia (stazione eretta, deambulazione);
  \item \textbf{medio gluteo}: cresta iliaca $\rightarrow$ grande trocantere; abduce ed extraruota;
    a punto fisso inclina il bacino;
  \item \textbf{piccolo gluteo}: linea glutea inferiore $\rightarrow$ grande trocantere;
    intraruota/abduce;
  \item \textbf{piriforme}: sacro/incisura ischiatica $\rightarrow$ grande trocantere; abduce ed
    extraruota, stabilizza l'anca;
  \item \textbf{otturatore esterno} (fossa otturatoria esterna) e \textbf{interno} (fossa interna)
    $\rightarrow$ grande trocantere; stabilizzatori (esterno extraruota, interno intraruota);
  \item \textbf{gemelli} superiore/inferiore $\rightarrow$ fossetta trocanterica; extraruotano,
    stabilizzano;
  \item \textbf{quadrato del femore}: tuberosità ischiatica $\rightarrow$ cresta intertrocanterica;
    extraruota, abduce.
\end{itemize}

\subsection{Muscoli della coscia}
\begin{itemize}
  \item \textbf{anteriori} (estensori): quadricipite femorale, sartorio;
  \item \textbf{mediali} (adduttori): gracile, pettineo, adduttore lungo/breve, grande adduttore;
  \item \textbf{posteriori} (flessori): bicipite femorale, semitendinoso, semimembranoso.
\end{itemize}
\begin{itemize}
  \item \textbf{quadricipite femorale}: 4 capi (retto femorale $\rightarrow$ spina iliaca ant.\ inf.;
    vasto laterale, mediale, intermedio $\rightarrow$ femore/linea aspra) $\rightarrow$ patella;
  \item \textbf{sartorio}: spina iliaca ant.\ sup.\ $\rightarrow$ condilo mediale della tibia; flette
    gamba e coscia, extraruota e abduce la coscia;
  \item \textbf{bicipite femorale}: capo lungo (tuberosità ischiatica) + breve (linea aspra)
    $\rightarrow$ condilo laterale della tibia e testa della fibula; flette ed extraruota la gamba;
  \item \textbf{semitendinoso} e \textbf{semimembranoso}: tuberosità ischiatica $\rightarrow$ condilo
    mediale della tibia; estendono la coscia, flettono e intraruotano la gamba;
  \item \textbf{gracile}: ramo ischio-pubico $\rightarrow$ condilo mediale della tibia; adduce,
    flette e intraruota;
  \item \textbf{pettineo}: cresta pettinea $\rightarrow$ linea pettinea del femore;
    extraruota/flette/adduce;
  \item \textbf{adduttore lungo} e \textbf{breve}: ramo superiore del pube $\rightarrow$ linea aspra;
    extraruotano/flettono/adducono;
  \item \textbf{grande adduttore}: tuberosità ischiatica $\rightarrow$ linea aspra e tubercolo del
    grande adduttore; extraruota e adduce.
\end{itemize}

\subsection{Triangolo femorale}
Regione triangolare della faccia anteriore prossimale della coscia. Limiti: legamento inguinale
(alto), sartorio (basso-laterale), adduttore lungo (basso-mediale); pavimento tra ileopsoas e
pettineo; rivestito dalla fascia lata (centralmente fascia cribrosa, attraversata dalla grande vena
safena). Contiene i linfonodi inguinali; sede dell'accesso all'arteria femorale (angioplastica) e
della linfoadenectomia inguinale; principali patologie: ernie crurali.

\subsection{Fossa poplitea}
Fossa posteriore al ginocchio, a rombo, con tessuto adiposo (linfonodi poplitei, nervo ischiatico,
vena e arteria poplitea). Limiti: bicipite femorale (alto-laterale), semitendinoso/semimembranoso
(alto-mediale), capi del gastrocnemio (in basso). Patologie: cisti poplitee, aneurismi dell'arteria/
vena poplitea, sindrome da intrappolamento dell'arteria poplitea.

\subsection{Muscoli della gamba}
\begin{itemize}
  \item \textbf{anteriori} (estensori): tibiale anteriore, estensore lungo delle dita, estensore
    lungo dell'alluce, peroniero anteriore;
  \item \textbf{laterali}: peroniero lungo e breve;
  \item \textbf{posteriori} (flessori): tricipite della sura, plantare, popliteo, flessore lungo
    delle dita, flessore lungo dell'alluce, tibiale posteriore.
\end{itemize}
Racchiusi nella fascia crurale (con retinacoli alla caviglia). A muovere la gamba è soprattutto la
coscia; eccezione = \textbf{popliteo} (intrarotatore della gamba flessa). Movimenti: plantarflessione,
dorsiflessione, flesso-estensione di gamba/dita, eversione/inversione del piede.
\begin{itemize}
  \item \textbf{tibiale anteriore}: condilo laterale/membrana interossea $\rightarrow$ I cuneiforme e
    I metatarsale; flette dorsalmente e ruota medialmente il piede;
  \item \textbf{estensore lungo dell'alluce}: fibula/membrana interossea $\rightarrow$ falange
    dell'alluce; lo estende, adduce il piede;
  \item \textbf{estensore lungo delle dita}: condilo laterale/fibula $\rightarrow$ ultime 4 falangi
    (4 tendini); le estende, ruota in fuori il piede;
  \item \textbf{peroniero anteriore}: fibula $\rightarrow$ V metatarsale; flette dorsalmente e ruota
    in fuori;
  \item \textbf{peroniero lungo}: testa della fibula $\rightarrow$ I metatarsale (pianta); flette
    plantarmente, abduce, ruota lateralmente; \textbf{breve}: $\rightarrow$ V metatarsale; abduce e
    ruota lateralmente;
  \item \textbf{tricipite della sura}: \textbf{gastrocnemio} (2 capi, femore $\rightarrow$ calcagno)
    + \textbf{soleo} (tibia/fibula $\rightarrow$ tendine calcaneale); flettono plantarmente, sollevano
    il corpo sulle punte;
  \item \textbf{plantare}: condilo laterale del femore $\rightarrow$ tuberosità calcaneale (col tendine
    d'Achille); coadiuva il tricipite;
  \item \textbf{popliteo}: dietro il ginocchio; flette e intraruota la gamba;
  \item \textbf{flessore lungo delle dita}: sotto la linea poplitea $\rightarrow$ ultime 4 falangi
    (pianta); le flette;
  \item \textbf{flessore lungo dell'alluce}: fibula/membrana interossea $\rightarrow$ falange
    dell'alluce; la flette;
  \item \textbf{tibiale posteriore}: tibia (sotto la linea poplitea) $\rightarrow$ navicolare e I
    cuneiforme; flette plantarmente e adduce il piede.
\end{itemize}
